\section{Convergence Diagnostics}
\label{sec:convergence}

Before using an ensemble for inference, one must verify that the Markov chain
has approximately converged to its stationary distribution.  We review three
diagnostics used in the G-track and developed in detail in G.4.

\subsection{Gelman-Rubin \texorpdfstring{$\rhat$}{R-hat}}

The \citet{gelman1992inference} potential scale reduction factor is computed
from $m \geq 2$ independent chains of length $n$.  Let $\bar{\psi}_{j\cdot}$
be the within-chain mean and $\bar{\psi}_{\cdot\cdot}$ the grand mean of
statistic $\psi$:

\begin{align}
  B &= \frac{n}{m-1} \sum_{j=1}^{m}(\bar{\psi}_{j\cdot} - \bar{\psi}_{\cdot\cdot})^2
      \quad \text{(between-chain variance)} \\
  W &= \frac{1}{m}\sum_{j=1}^{m} s_j^2
      \quad \text{(within-chain variance)} \\
  \hat{V} &= \frac{n-1}{n} W + \frac{m+1}{mn} B \\
  \rhat &= \sqrt{\hat{V}/W}.
\end{align}

A value $\rhat < 1.1$ is conventionally accepted as indicating convergence
for redistricting's discrete integer outcomes.\footnote{Vehtari et al.\ (2021)
recommend $\rhat < 1.01$ for continuous parameters; 1.1 is the appropriate
threshold for the discrete seat-count and plan-assignment statistics used
here.}
For redistricting ensembles, we compute $\rhat$ over partisan seats and
Polsby-Popper means.

\subsection{Effective Sample Size (ESS)}

The effective sample size corrects for autocorrelation within a chain.  For
a chain of length $N$ with autocorrelation $\rho_t$ at lag $t$:

\begin{equation}
  \ess = \frac{N}{1 + 2\sum_{t=1}^{\infty} \rho_t}.
  \label{eq:ess}
\end{equation}

An \ess{} of at least 400 per statistic is required for reliable 95\%
interval estimates.  \recom{} ensembles typically achieve higher \ess{} than
flip-chain ensembles of the same length, because the recombination move
decorrelates faster.

\subsection{Hamming Autocorrelation}

District-level statistics like seat share depend on which tracts belong to
which district, not just on continuous quantities.  The \emph{Hamming
distance} between two plans $P$ and $P'$ is:

\begin{equation}
  d_H(P, P') = \frac{1}{n} \#\{i : \text{district}(t_i, P) \neq
               \text{district}(t_i, P')\},
  \label{eq:hamming}
\end{equation}

where $t_i$ ranges over census tracts and we match districts by population
overlap.  The \emph{Hamming autocorrelation} at lag $\ell$ is:

\begin{equation}
  \rho_H(\ell) = \frac{\mathrm{Cov}[d_H(P_t, P_{\text{ref}}),
                  d_H(P_{t+\ell}, P_{\text{ref}})]}
                 {\mathrm{Var}[d_H(P_t, P_{\text{ref}})]},
  \label{eq:hamming-autocorr}
\end{equation}

where $P_{\text{ref}}$ is a fixed reference plan.  Rapid decay of
$\rho_H(\ell)$ to zero indicates that the chain is exploring the plan space
efficiently.  G.4 reports Hamming autocorrelation plots for all six states.

\subsection{Diagnostic Standards for the G-Track}

The following standards apply across all G-track papers:

\begin{center}
\begin{tabular}{lll}
\toprule
Diagnostic & Threshold & Consequence of failure \\
\midrule
$\rhat$ & $< 1.1$ & Results flagged as preliminary \\
\ess    & $> 400$  & Wider uncertainty intervals used \\
$\rho_H(10)$ & $< 0.10$ & Chain considered mixed \\
\bottomrule
\end{tabular}
\end{center}

For published ensembles (Herschlag 2020, DeFord 2021) we reproduce the
authors' reported convergence diagnostics.  Where these are unavailable,
we apply conservative uncertainty adjustments to reported percentiles.
